\section*{Chapitre \ref{ch:g9:weight-vs-mass} --- Poids et masse~: mesurer les forces}

\begin{solution}{exo:g9:weight-vs-mass:1}
Le \omterm{def:g9:weight-vs-mass:newton}{newton} (\unit{N}), mesuré au \omterm{def:g9:weight-vs-mass:dynamometer}{dynamomètre}. Repères~: un \omterm{def:g9:weight-vs-mass:newton}{newton} ---
le \omterm{def:g4:weight-and-mass:weight}{poids} d'une petite pomme~; cent \omterm{def:g9:weight-vs-mass:newton}{newtons} --- une poignée de main
ferme (ou le \omterm{def:g4:weight-and-mass:weight}{poids} d'une valise de \qty{10}{kg}).
\end{solution}

\begin{solution}{exo:g9:weight-vs-mass:2}
$P = m \times g$~: le \omterm{def:g4:weight-and-mass:weight}{poids} en \omterm{def:g9:weight-vs-mass:newton}{newtons}, la \omterm{def:g3:measuring-mass:mass}{masse} en \omterm{def:g3:measuring-mass:units}{kilogrammes}, et
$g$ --- l'\omterm{prop:g9:weight-vs-mass:formula}{intensité de la pesanteur} du lieu --- les \omterm{def:g9:weight-vs-mass:newton}{newtons}
d'attraction que chaque \omterm{def:g3:measuring-mass:units}{kilogramme} y reçoit (\qty{9.8}{N/kg} sur
Terre).
\end{solution}

\begin{solution}{exo:g9:weight-vs-mass:3}
Chien~: $25 \times 9.8 = \qty{245}{N}$. Pain~:
$0.7 \times 9.8 \approx \qty{6.9}{N}$. Un lecteur de \qty{55}{kg}~:
environ \qty{539}{N}.
\end{solution}

\begin{solution}{exo:g9:weight-vs-mass:4}
$m = 14.7 \div 9.8 = \qty{1.5}{kg}$.
\end{solution}

\begin{solution}{exo:g9:weight-vs-mass:5}
La pente est le $g$ local --- des \omterm{def:g9:weight-vs-mass:newton}{newtons} par \omterm{def:g3:measuring-mass:units}{kilogramme}. Sur la \omterm{def:g2:moon-in-the-sky:moon}{Lune},
la droite reste droite et passe toujours par l'origine, mais
s'affaisse à une pente douce de $1.6$~: même proportionnalité, monde
plus faible.
\end{solution}

\begin{solution}{exo:g9:weight-vs-mass:6}
La colonne des \omterm{def:g3:measuring-mass:mass}{masses} ne change jamais --- \qty{10}{kg} de matière
voyagent intacts. Les colonnes de $g$ et du \omterm{def:g4:weight-and-mass:weight}{poids} changent avec le
monde~: le \omterm{def:g4:weight-and-mass:weight}{poids} est le péage local, $m \times g$.
\end{solution}

\begin{solution}{exo:g9:weight-vs-mass:7}
La \omterm{def:g3:levers-and-scales:scale}{balance} à fléau compare $m_1 g$ à $m_2 g$~: $g$ multiplie les deux
plateaux et se simplifie --- la \omterm{def:g3:measuring-mass:mass}{masse}, véridiquement, partout. Le
pèse-personne mesure $P = m g_{\text{ici}}$ mais divise par
$g_{\text{Terre}}$ pour afficher des \omterm{def:g3:measuring-mass:units}{kilogrammes}~: sur la \omterm{def:g2:moon-in-the-sky:moon}{Lune}, il
affiche $m \times 1.6/9.8$ --- le sixième de la vérité.
\end{solution}

\begin{solution}{exo:g9:weight-vs-mass:8}
\omterm{def:g4:weight-and-mass:weight}{Poids}~: \qty{19.6}{N} vers le bas --- une flèche d'environ
\qty{3.9}{cm} à cinq \omterm{def:g9:weight-vs-mass:newton}{newtons} par \omterm{def:g2:measuring-length:units}{centimètre}, partant du livre vers le
bas. Soutien~: une flèche égale de \qty{3.9}{cm} vers le haut.
Longueurs égales, sens opposés~: budget équilibré, livre immobile.
\end{solution}

\begin{solution}{exo:g9:weight-vs-mass:9}
L'avion doit porter, soulever et freiner la \emph{matière} du bagage
--- sa \omterm{def:g3:measuring-mass:mass}{masse} --- et \qty{23}{kg} de matière font \qty{23}{kg} sur
n'importe quel monde. Le comptoir lunaire a besoin d'une \omterm{def:g3:levers-and-scales:scale}{balance} à
fléau avec des \omterm{def:g3:measuring-mass:mass}{masses} marquées~: la seule qui survive au changement de
$g$.
\end{solution}

\begin{solution}{exo:g9:weight-vs-mass:10}
$m = 111 \div 3.7 = \qty{30}{kg}$~; sur Terre, il pèse
$30 \times 9.8 = \qty{294}{N}$.
\end{solution}

\begin{solution}{exo:g9:weight-vs-mass:11}
Il indique trop~: suspendu, le ressort porte aussi son propre \omterm{def:g4:weight-and-mass:weight}{poids}
(et celui de son crochet), que le zéro fait à plat n'a jamais compté
--- l'instrument ajoute un excédent constant. Mets-le à zéro suspendu,
et l'excédent est retranché avant même l'arrivée de la charge.
\end{solution}

\begin{solution}{exo:g9:weight-vs-mass:12}
Matériel~: un \omterm{def:g9:weight-vs-mass:dynamometer}{dynamomètre}, un jeu de \omterm{def:g3:measuring-mass:mass}{masses} marquées ($1$ à
\qty{5}{kg}), un crochet. Suspendre chaque \omterm{def:g3:measuring-mass:mass}{masse}, noter les couples
($m$, $P$), tracer $P$ en fonction de $m$~: les points s'alignent en
passant par l'origine, et la pente --- la montée en \omterm{def:g9:weight-vs-mass:newton}{newtons} sur
l'avancée en \omterm{def:g3:measuring-mass:units}{kilogrammes} --- vaut $g \approx \qty{9.8}{N/kg}$. Sur la
\omterm{def:g2:moon-in-the-sky:moon}{Lune}~: protocole identique, alignement identique, pente nouvelle de
$1.6$ --- la méthode mesure le monde sur lequel elle se tient~; seul
le nombre est local.
\end{solution}

\begin{solution}{pb:g9:weight-vs-mass:1}
\textbf{1.} Terre $70 \times 9.8 = \qty{686}{N}$~; \omterm{def:g2:moon-in-the-sky:moon}{Lune}
\qty{112}{N}~; Mars \qty{259}{N}~; pont de Jupiter \qty{1736}{N}.
\textbf{2.} Correct sur Terre~: \qty{70}{kg}. \omterm{def:g2:moon-in-the-sky:moon}{Lune}~:
$112 \div 9.8 \approx \qty{11.4}{kg}$~; Mars~:
$259 \div 9.8 \approx \qty{26.4}{kg}$~; Jupiter~:
$1736 \div 9.8 \approx \qty{177}{kg}$ --- la flatterie et la calomnie
par formule.
\textbf{3.} \qty{70}{kg} à tous les ports~: la comparaison de la
\omterm{def:g3:levers-and-scales:scale}{balance} à fléau élimine le $g$ local.
\textbf{4.} Seulement le pont de Jupiter~: \qty{1736}{N} approche la
limite mais reste en dessous --- aucun port ne dépasse \qty{2000}{N}
pour ce voyageur (un collègue plus lourd, de \qty{90}{kg},
l'enfreindrait là-bas~: $90 \times 24.8 = \qty{2232}{N}$).
\textbf{5.} Départ~: $40 \times 9.8 = \qty{392}{N}$~; arrivée~:
$40 \times 3.7 = \qty{148}{N}$.
\textbf{6.} $m = 200/g$~: Terre \qty{20.4}{kg}~; \omterm{def:g2:moon-in-the-sky:moon}{Lune} \qty{125}{kg}~;
Mars \qty{54}{kg}~; Jupiter \qty{8.1}{kg} --- les \omterm{def:g9:weight-vs-mass:newton}{newtons} fixes du
syndicat achètent des \omterm{def:g3:measuring-mass:mass}{masses} très différentes.
\textbf{7.} $m = 150 \div 1.6 = \qty{93.75}{kg}$ --- environ
\qty{94}{kg} de caisse pour la douane.
\textbf{8.} $m = 294 \div 9.8 = \qty{30}{kg}$~; \omterm{def:g4:weight-and-mass:weight}{poids} lunaire
$30 \times 1.6 = \qty{48}{N}$.
\textbf{9.} Les muscles et la \omterm{def:g3:measuring-mass:mass}{masse} sont rentrés inchangés~; seul
l'adversaire a changé. Sur la \omterm{def:g2:moon-in-the-sky:moon}{Lune}, chaque \omterm{def:g3:measuring-mass:units}{kilogramme} coûtait
\qty{1.6}{N} à soulever~; la maison facture de nouveau \qty{9.8}{N}.
Les gains étaient la remise de la \omterm{def:g2:moon-in-the-sky:moon}{Lune}, non la croissance du corps ---
$P = m g$, avec $m$ fidèle et $g$ volage.
\textbf{10.} $0.2 \times 24.8 \approx \qty{5}{N}$ --- une enclume,
non~: sur Terre le fouet pesait \qty{2}{N}, si bien que battre sur le
pont est deux fois et demie plus lourd, fatigant pour un poignet mais
guère un travail de forge. Réclamation~: partiellement retenue.
\textbf{11.} Un \omterm{def:g9:weight-vs-mass:newton}{newton} de chocolat vaut $m = 1/g$~: Terre environ
\qty{102}{g}~; \omterm{def:g2:moon-in-the-sky:moon}{Lune} \qty{625}{g}~; Jupiter \qty{40}{g}. Achète tes
\omterm{def:g9:weight-vs-mass:newton}{newtons} sur la \omterm{def:g2:moon-in-the-sky:moon}{Lune}.
\textbf{12.} Par exemple~: «~La \omterm{def:g3:measuring-mass:mass}{masse} voyage avec toi~; le \omterm{def:g4:weight-and-mass:weight}{poids} est
le péage que chaque monde prélève dessus --- emporte des \omterm{def:g3:measuring-mass:units}{kilogrammes},
et laisse les ports facturer leurs \omterm{def:g9:weight-vs-mass:newton}{newtons}.~»
\end{solution}
